\documentclass[12pt]{article}
\usepackage[utf8]{inputenc}
\usepackage[T1]{fontenc}
\usepackage[margin=1in]{geometry}
\usepackage{setspace}
\usepackage{xcolor}
\usepackage[colorlinks=true, urlcolor=blue!60!black, linkcolor=blue!60!black]{hyperref}
\setlength{\parindent}{2em}
\setlength{\parskip}{1.25em}
\renewcommand{\baselinestretch}{1.0}
\title{\textit{Counterman v. Colorado}, 600 U.S. 66 (2023)}
\author{Case Brief}
\date{}
\begin{document}
\maketitle

% Case Brief: Counterman v. Colorado, 600 U.S. 66 (2023)
% Chapter: Free Speech & Press
% Prepared for Seth J. Chandler's Constitutional Law Briefing Book

\section{Case Metadata}

\begin{itemize}
  \item \textbf{Case Name:} \href{https://supreme.justia.com/cases/federal/us/600/22-138/}{\textit{Counterman v. Colorado}} (Billy Raymond Counterman, Petitioner v. Colorado)
  \item \textbf{Citation:} 600 U.S. 66 (2023); Docket No. 22-138 (\href{https://www.supremecourt.gov/opinions/22pdf/600us1r51_g3bi.pdf}{Official Supreme Court Slip/Page Proof PDF})
  \item \textbf{Dates:} Certiorari Granted January 13, 2023; Argued April 19, 2023; Decided June 27, 2023
  \item \textbf{Court Below:} Court of Appeals of Colorado, 497 P.3d 1039 (Colo. App. 2021) (Colorado Supreme Court denied review)
  \item \textbf{Disposition:} Vacated and remanded (7--2)
  \item \textbf{Vote and Alignment:}
  \begin{itemize}
    \item \textbf{Majority (5):} Justice Elena Kagan (author), joined by Chief Justice John G. Roberts, Jr., and Justices Samuel A. Alito, Jr., Brett M. Kavanaugh, and Ketanji Brown Jackson.
    \item \textbf{Concurrence in Part and in the Judgment (2):} Justice Sonia Sotomayor (author), joined by Justice Neil M. Gorsuch as to Parts I, II, III-A, and III-B.
    \item \textbf{Dissents (2):} Justice Clarence Thomas (separate dissent, joining Justice Barrett in full); Justice Amy Coney Barrett (dissenting opinion, joined by Justice Thomas).
  \end{itemize}
  \item \textbf{Constitutional Provision \& Statute:} First Amendment (Free Speech Clause, incorporated via Fourteenth Amendment); Colo. Rev. Stat. \S~18-3-602(1)(c) (2022).
\end{itemize}

\section{Facts and Procedural History}

From 2014 to 2016, petitioner Billy Raymond Counterman sent hundreds of unsolicited Facebook messages to C.W. (local Colorado singer-songwriter Coles Whalen, who participated as \textit{amicus curiae} at the Supreme Court) (\href{https://supreme.justia.com/cases/federal/us/600/22-138/}{600 U.S. at 70}). The two had never met, and Whalen never replied. Each time she blocked him, Counterman created a new Facebook account and resumed contacting her. Some messages were facially mundane (``Good morning sweetheart''), others implied physical surveillance (``Was that you in the white Jeep?''; ``a couple [of] physical sightings''), and several expressed anger and envisaged violent harm: ``Fuck off permanently''; ``Staying in cyber life is going to kill you''; and ``You're not being good for human relations. Die.'' (\href{https://www.supremecourt.gov/opinions/22pdf/600us1r51_g3bi.pdf}{600 U.S. at 70}).

The messages terrified Whalen and upended her life: she suffered severe anxiety and insomnia, stopped walking alone, declined social engagements, and canceled musical performances at financial cost before contacting authorities (\href{https://supreme.justia.com/cases/federal/us/600/22-138/}{600 U.S. at 70}). Colorado charged Counterman under its stalking statute, Colo. Rev. Stat. \S~18-3-602(1)(c), which criminalizes ``[r]epeatedly . . . mak[ing] any form of communication with another person'' in a manner that would cause a reasonable person serious emotional distress and actually does so. Because Colorado lacked evidence of physical following, its prosecution relied solely on Counterman's Facebook messages (\href{https://supreme.justia.com/cases/federal/us/600/22-138/}{600 U.S. at 70 \& n.1}).

Counterman moved to dismiss on First Amendment grounds, arguing his messages were not ``true threats.'' Applying Colorado's purely objective ``reasonable person'' test (\textit{People v. Cross}, 127 P.3d 71, 76 (Colo. 2006); \textit{In re R.D.}, 464 P.3d 717, 731 n.21 (Colo. 2020)), the trial court held the State did not need to prove subjective intent to threaten; a jury convicted him, and he received a four-and-a-half-year prison sentence. The Colorado Court of Appeals affirmed (497 P.3d 1039 (2021)), and the Colorado Supreme Court denied review. The U.S. Supreme Court granted certiorari to resolve a deep circuit and state-court split left open after \href{https://supreme.justia.com/cases/federal/us/575/723/}{\textit{Elonis v. United States}, 575 U.S. 723 (2015)} (\href{https://supreme.justia.com/cases/federal/us/600/22-138/}{600 U.S. at 71--72}).

\section{Issue(s) Presented}

\begin{enumerate}
  \item Whether the First Amendment requires the government in a criminal ``true threats'' prosecution to prove that the speaker had some subjective understanding of the threatening nature of his statements, rather than relying solely on an objective reasonable-person standard.
  \item If a subjective \textit{mens rea} is constitutionally required, whether recklessness suffices or whether the First Amendment demands purpose or knowledge.
\end{enumerate}

\section{Holding}

\textbf{Yes to subjective \textit{mens rea}; recklessness is sufficient; judgment vacated and remanded (7--2).} Writing for the Court, Justice Kagan held that although ``true threats'' of violence are historically unprotected speech defined objectively by what the statement conveys, the First Amendment requires the State in true-threats prosecutions to prove the defendant had a subjective mental state of at least \textbf{recklessness}---meaning the defendant ``consciously disregarded a substantial risk that his communications would be viewed as threatening violence'' (\href{https://www.supremecourt.gov/opinions/22pdf/600us1r51_g3bi.pdf}{600 U.S. at 69}). Because Colorado prosecuted Counterman under an objective negligence standard without proving any subjective awareness of his messages' threatening character, his conviction violated the First Amendment (\href{https://supreme.justia.com/cases/federal/us/600/22-138/}{600 U.S. at 82--83}).

\section{Reasoning}

\subsection{Why a Subjective Mental State Is Required: ``Strategic Protection'' Against Chill}

Justice Kagan began from settled categorical doctrine: ``true threats''---serious expressions conveying that a speaker means to commit unlawful violence (\textit{Virginia v. Black}, 538 U.S. 343, 359 (2003))---are historically unprotected speech, distinct from jests or political hyperbole (\textit{Watts v. United States}, 394 U.S. 705, 708 (1969)). Crucially, the Court treated the \textit{existence} of a threat as objective, turning on ``what the statement conveys'' to the recipient rather than the speaker's mindset (\href{https://supreme.justia.com/cases/federal/us/575/723/}{\textit{Elonis}, 575 U.S. at 733}).

Nevertheless, the First Amendment frequently provides ``strategic protection''---a prophylactic ``buffer zone'' shielding some unprotected speech to prevent a \textbf{chilling effect} on protected expression (\textit{Gertz v. Robert Welch, Inc.}, 418 U.S. 323, 342 (1974); \textit{Speiser v. Randall}, 357 U.S. 513, 526 (1958)). Fear of misjudging the legal line, judicial error, or litigation expense induces speakers to self-censor lawful speech. Kagan surveyed three analogous unprotected categories where the Court requires subjective fault:
\begin{itemize}
  \item \textbf{Defamation:} \textit{New York Times Co. v. Sullivan}, 376 U.S. 254, 280 (1964), requires ``actual malice'' (knowledge or reckless disregard of falsity) for public-figure defamation and criminal libel (\textit{Garrison v. Louisiana}, 379 U.S. 64, 74 (1964)).
  \item \textbf{Incitement:} \textit{Brandenburg v. Ohio}, 395 U.S. 444, 447 (1969), and \textit{Hess v. Indiana}, 414 U.S. 105, 109 (1973), require specific intent to produce imminent lawless action.
  \item \textbf{Obscenity:} \textit{Hamling v. United States}, 418 U.S. 87, 122--123 (1974), and \textit{Smith v. California}, 361 U.S. 147, 151--152 (1959), demand scienter---awareness of the material's character and nature.
\end{itemize}
Applying the same logic to true threats, an objective standard (effectively negligence) forces speakers to give threats ``a wide berth,'' chilling heated, non-threatening discourse (\href{https://supreme.justia.com/cases/federal/us/600/22-138/}{600 U.S. at 75--78}).

\subsection{Why Recklessness---Not Purpose or Knowledge---Strikes the Right Balance}

Descending the \textit{mens rea} hierarchy (purpose, knowledge, recklessness, negligence), Kagan selected \textbf{recklessness}---where a speaker is aware ``that others could regard his statements as'' threatening violence and ``delivers them anyway'' (\href{https://supreme.justia.com/cases/federal/us/575/723/}{\textit{Elonis}, 575 U.S. at 746} (Alito, J., concurring in part and dissenting in part))---for three reasons (\href{https://supreme.justia.com/cases/federal/us/600/22-138/}{600 U.S. at 78--82}):
\begin{enumerate}
  \item \textbf{Protecting Victims from Profound Harm:} Elevating the floor to purpose or knowledge would severely impede prosecutions of morally culpable speakers (especially absent direct evidence of intent) with diminishing returns for protected speech.
  \item \textbf{Alignment with \textit{Sullivan} Defamation:} Recklessness has provided ample ``breathing space'' in public-figure defamation for over half a century, and true threats deserve no greater constitutional insulation than defamation.
  \item \textbf{Distinguishing Incitement:} Incitement requires specific intent because it sits ``a hair's-breadth away'' from core political advocacy (\textit{Brandenburg}; \textit{NAACP v. Claiborne Hardware Co.}, 458 U.S. 886 (1982)). Speech near the true-threats line is typically less central to First Amendment theory and less vulnerable to political suppression.
\end{enumerate}

\section{Concurring and Dissenting Opinions}

\begin{itemize}
  \item \textbf{Justice Sotomayor (Concurring in Part and in the Judgment, joined by Justice Gorsuch as to Parts I, II, III-A, and III-B):} Sotomayor agreed that subjective \textit{mens rea} is required and that recklessness suffices for Counterman's \textit{stalking} prosecution, because repeated, unwanted, targeted contact raises fewer First Amendment concerns than isolated pure-speech utterances. She criticized the majority for reaching pure true threats generally without party briefing on recklessness. Grounded in \textit{Virginia v. Black} and historical threat statutes, she argued (alone in Parts III-C, III-D, and IV) that pure true threats require \textbf{purpose or knowledge}; lowering the bar to recklessness risks overcriminalizing online hyperbole, rap lyrics, and political protest---effectively letting prosecutors evade \textit{Brandenburg}'s intent rule by recharging incitement as threats (\href{https://supreme.justia.com/cases/federal/us/600/22-138/}{600 U.S. at 83--105}).
  \item \textbf{Justice Thomas (Dissenting):} Joining Justice Barrett's dissent in full, Thomas wrote separately to denounce the majority's ``surprising and misplaced reliance'' on \textit{New York Times Co. v. Sullivan}, reiterating his view that \textit{Sullivan} is a ``policy-driven'' departure from Founding-era common law that the Court should reconsider rather than extend to true threats (\href{https://supreme.justia.com/cases/federal/us/600/22-138/}{600 U.S. at 105--106}).
  \item \textbf{Justice Barrett (Dissenting, joined by Justice Thomas):} Barrett argued that an \textbf{objective reasonable-person standard} is constitutionally sufficient. Most unprotected categories (fighting words, false commercial speech, private-figure defamation, and obscenity's community-standards test) use objective standards, and contextual review under \textit{Watts} and \textit{Black} already prevents chilling protected hyperbole. Neither history nor precedent supports a recklessness floor, which Barrett dismissed as a legislative ``Goldilocks judgment.'' She warned that because \textit{Counterman} is a constitutional ruling, it threatens to obstruct civil protective orders, civil enforcement suits (e.g., 18 U.S.C. \S~248), and school or workplace discipline against delusional or evasive threat-makers (\href{https://supreme.justia.com/cases/federal/us/600/22-138/}{600 U.S. at 106--122}).
\end{itemize}

\section{Key Quotations}

\begin{enumerate}
  \item \textbf{The Constitutional Floor of Recklessness (Kagan, J., \href{https://www.supremecourt.gov/opinions/22pdf/600us1r51_g3bi.pdf}{600 U.S. at 69}):}
  \begin{quote}
    ``We hold that it does, but that a mental state of recklessness is sufficient. The State must show that the defendant consciously disregarded a substantial risk that his communications would be viewed as threatening violence. The State need not prove any more demanding form of subjective intent to threaten another.''
  \end{quote}

  \item \textbf{Prophylactic Protection Against Chilling Effects (Kagan, J., \href{https://www.supremecourt.gov/opinions/22pdf/600us1r51_g3bi.pdf}{600 U.S. at 75}):}
  \begin{quote}
    ``Yet the First Amendment may still demand a subjective mental-state requirement shielding some true threats from liability. The reason relates to what is often called a chilling effect. Prohibitions on speech have the potential to chill, or deter, speech outside their boundaries.''
  \end{quote}

  \item \textbf{Limiting Recklessness to Stalking (Sotomayor, J., concurring in part and in the judgment, \href{https://supreme.justia.com/cases/federal/us/600/22-138/}{600 U.S. at 85}):}
  \begin{quote}
    ``There is simply no need to reach out in this stalking case to determine whether anything more than recklessness is needed for punishing true threats generally.''
  \end{quote}

  \item \textbf{Critique of the Majority's ``Goldilocks Judgment'' (Barrett, J., dissenting, \href{https://supreme.justia.com/cases/federal/us/600/22-138/}{600 U.S. at 119}):}
  \begin{quote}
    ``The reality is that recklessness is not grounded in law, but in a Goldilocks judgment: Recklessness is not too much, not too little, but instead `just right.'''
  \end{quote}
\end{enumerate}

\section{Doctrinal Significance and Modern Legacy (Through 2026)}

\textit{Counterman} resolved the constitutional question left open by \href{https://supreme.justia.com/cases/federal/us/575/723/}{\textit{Elonis}} and unified federal and state true-threats prosecutions under a single two-part constitutional test: (1) an objective true threat in context, plus (2) subjective recklessness. Between 2023 and 2026, four major developments have defined its reach:
\begin{itemize}
  \item \textbf{Content-Neutral Stalking Carve-Out:} Following Justice Sotomayor's concurrence, state high courts have cabined \textit{Counterman} to content-based threat prosecutions. In \href{https://law.justia.com/cases/colorado/supreme-court/2025/24sa226.html}{\textit{People v. Crawford} (Colo. 2025)}, the Colorado Supreme Court held that \textit{Counterman} is ``simply inapplicable'' to stalking prosecutions resting on non-content conduct or repeated unwanted contact rather than threatening speech content.
  \item \textbf{Reaffirming \textit{Brandenburg} for Political Protest:} When denying certiorari in \textit{Mckesson v. Doe} (2024), Justice Sotomayor underscored \textit{Counterman}'s explicit reaffirmation that negligence cannot ground liability for political speech and that incitement demands specific intent---guiding remand litigation through 2026 (\href{https://www.aclu.org/press-releases/fifth-circuit-rules-against-civil-rights-activist-deray-mckesson-in-long-running-protest-rights-case}{ACLU Case Summary, Mar. 2026}).
  \item \textbf{Jury Instructions and Survival of Pre-\textit{Counterman} Statutes:} Through 2026, federal and state appellate courts have largely rejected facial invalidation of objective-worded threat statutes, curing constitutional defects via recklessness jury instructions and applying harmless-error review (e.g., \href{https://www.supremecourt.gov/DocketPDF/26/26-411/425007/20260921123229913_Calloway%20Petition%20September%2021%202026%20EFile.pdf}{\textit{State v. Calloway}, 585 P.3d 1270 (Wash. 2026)}; \href{https://www.supremecourt.gov/DocketPDF/26/26A248/419828/20260821173446043_Thomas%20Woollacott%20Application%20for%20Ext%20to%20File%20Petition.pdf}{\textit{United States v. Woollacott} (6th Cir. 2026)}).
  \item \textbf{Contested Civil Boundary:} Empirical studies through 2026 report that while most state courts have declined to extend \textit{Counterman}'s recklessness floor to civil protective orders, a growing minority have done so---leaving Justice Barrett's warning about civil injunctions an active fault line (\href{https://scholarlycommons.law.northwestern.edu/cgi/viewcontent.cgi?article=7800&context=jclc}{Birnholtz \& Kugler, 116 J. Crim. L. \& Criminology 362 (2026)}; \href{https://www.californialawreview.org/print/stalking-counterman}{Lakier \& Douek, 113 Cal. L. Rev. 143 (2025)}).
\end{itemize}

\section{Pedagogical Critiques and Exam Hypotheticals}

\subsection{Key Pedagogical Critiques}
\begin{itemize}
  \item \textbf{The Delusional Stalker Paradox:} As Justice Barrett noted, genuinely delusional stalkers (like Counterman, who appeared to believe he was in an ongoing dialogue with Whalen) may pose the gravest danger to victims while lacking subjective awareness that their words are frightening.
  \item \textbf{Doctrinal Tension with History-and-Tradition:} Unlike \textit{Bruen} or \textit{Stevens}, Kagan's majority openly embraced pragmatic interest-balancing (``having much of what is important on both sides of the scale'') to craft a prophylactic buffer zone, drawing fire from both textualist-originalists (Thomas/Barrett) and speech-absolutists (Sotomayor/Gorsuch).
\end{itemize}

\subsection{Exam Hypotheticals (Including ``Opposite Side'' Contrast)}

\subsubsection{Hypothetical 1 (Online Hyperbole vs. Reckless Threat)}
During a heated livestreamed debate over a local zoning ordinance, a resident posts in the city council's public chat: ``Councilman Rivera is selling out our neighborhood---someone ought to burn his house down tonight.'' Prosecutors charge the resident under a state felony threat statute instructing the jury to convict if a ``reasonable viewer would regard the comment as a serious threat of violence.''

\textbf{Analysis:} The conviction violates \href{https://supreme.justia.com/cases/federal/us/600/22-138/}{\textit{Counterman}} because the jury instruction uses an objective negligence standard. On remand, two distinct hurdles arise: (1) under \textit{Watts}, was the statement objectively a ``true threat'' or protected political hyperbole? (2) If framed as \textit{incitement} (urging ``someone'' to act), \textit{Counterman} explicitly preserves \textit{Brandenburg}'s higher specific-intent standard, though prosecutors may try to bypass \textit{Brandenburg} by charging it as a reckless threat---precisely the danger Justice Sotomayor warned against.

\subsubsection{Hypothetical 2 (``Opposite Side'' Contrast---Protected Delusional Speech vs. Prosecutable Stalking)}
\begin{itemize}
  \item \textbf{Scenario A (Protected Under \textit{Counterman}):} After a single encounter at a book signing, an author receives one rambling email from a reader experiencing paranoid delusions: ``The universe requires your silence; your final chapter ends Friday.'' Medical experts establish the reader genuinely believed he was transmitting a harmless spiritual prophecy and never perceived any risk of causing fear. Under \textit{Counterman}, a criminal true-threats conviction fails for lack of subjective recklessness.
  \item \textbf{Scenario B (``Opposite Side''---Constitutionally Prosecutable):} Same reader sends 150 messages over six months despite the author blocking ten accounts and sending a cease-and-desist letter warning, ``Your messages terrify me; stop contacting me.'' Even if the reader claims he meant no harm, continuing to send ominous messages after explicit notice easily supports a jury finding of \textbf{recklessness} under \textit{Counterman}---and if prosecuted purely for repeated, unwanted post-notice contact rather than message content, \href{https://law.justia.com/cases/colorado/supreme-court/2025/24sa226.html}{\textit{Crawford}} indicates \textit{Counterman}'s true-threats test may not even be triggered.
\end{itemize}

\end{document}
